% !TEX root = task-2-11-3.tex
% Editable task bank. Append new BankTask blocks at the end of the document.
\documentclass[12pt,a4paper]{article}
\usepackage{physics-textbook}
\usepackage{graphicx,needspace}
\usetikzlibrary{decorations.pathmorphing}
\newif\ifShowSolutions
\ifdefined\TutorEdition
  \ShowSolutionstrue
\else
  \ShowSolutionsfalse % Change to \ShowSolutionstrue to display solutions.
\fi
\newcounter{banktask}
\newsavebox{\BankTaskBox}
\newcommand{\BankTaskContent}[3]{%
  \noindent\textbf{\thebanktask.}\quad #2\par
  \ifShowSolutions
    \Needspace{50mm}% Keep the solution heading with its opening text or diagram.
    #3\par
    \ifdefined\TutorEdition
      \ifcsname TutorSource#1\endcsname
        \nopagebreak\smallskip{\small\csname TutorSource#1\endcsname\par}%
      \fi
    \fi
  \fi
}
\newcommand{\BankTask}[3]{%
  \refstepcounter{banktask}\par\medskip
  \sbox{\BankTaskBox}{\begin{minipage}{\linewidth}%
    \BankTaskContent{#1}{#2}{#3}\end{minipage}}%
  \ifdim\dimexpr\ht\BankTaskBox+\dp\BankTaskBox\relax>\dimexpr\textheight-20mm\relax
    \Needspace{35mm}\BankTaskContent{#1}{#2}{#3}%
  \else
    \Needspace{\dimexpr\ht\BankTaskBox+\dp\BankTaskBox\relax}%
    \noindent\usebox{\BankTaskBox}\par
  \fi
  \medskip
}
\begin{document}
\BookSubsection{2.11.3}{Работа, совершаемая над газом}


% Source BA6eFF: https://ege.fipi.ru/bank/questions.php?search=1&proj=BA1F39653304A5B041B656915DC36B38&qid=BA6eFF
\BankTask{BA6eFF}{
Двум моль газа сообщили количество теплоты, равное \(1700\ \text{Дж}\), при этом внешние силы совершили над газом работу, равную \(500\ \text{Дж}\). На сколько увеличилась внутренняя энергия газа (в $\text{Дж}$)?
}{
\smallskip\noindent\textbf{Решение от GPT}\par
По первому закону термодинамики
\[
\Delta U=Q+A_{\text{внеш}}=1700+500=2200\ \text{Дж}.
\]
Ответ: \(2200\ \text{Дж}\).
}

% Source 5AD1F4: https://ege.fipi.ru/bank/questions.php?search=1&proj=BA1F39653304A5B041B656915DC36B38&qid=5AD1F4
\BankTask{5AD1F4}{
Газ в сосуде сжали, совершив работу, равную \(500\,\text{Дж}\). Внутренняя энергия газа при этом увеличилась на \(350\,\text{Дж}\). Какое количество теплоты отдал газ окружающей среде (в $\text{Дж}$)?
}{
\smallskip\noindent\textbf{Решение от GPT}\par
По первому закону термодинамики \(\Delta U=Q+A_{\text{внеш}}\), где \(Q\) — теплота, полученная газом. Тогда \(Q=350-500=-150\,\text{Дж}\). Газ отдал окружающей среде \(Q_{\text{отд}}=150\,\text{Дж}\).

Ответ: \(150\,\text{Дж}\).
}

\end{document}
